\cleardoublepage
\thispagestyle{empty}

\begin{center}
  \bfseries
  {\selectlanguage{german}\"Ubersicht}
\end{center}
\normalfont

\vspace{0.5cm}

{\selectlanguage{german}

Diese Arbeit untersucht die automatische Extraktion von Schl\"ussel-Wert-Paaren (Key-Value Pairs, KVPs) aus Gesch\"aftsdokumenten mit generativen und layoutbewussten diskriminativen Verfahren. Das diskriminative System kombiniert einen LayoutLMv3-Base-Encoder, einen Token-Klassifikator und einen Span-basierten Verkn\"upfer. Es verwendet extrahierten Text und Begrenzungsrahmen mit einer konstanten leeren visuellen Eingabe und erh\"alt daher keine dokumentspezifischen Bildinformationen.

\textbf{Hintergrund:} Die automatische Verarbeitung von Gesch\"aftsdokumenten stellt eine erhebliche Herausforderung im Bereich des Dokumentenverstehens dar. Traditionelle Ans\"atze basieren auf Vorlagen- oder regelbasierten Systemen, die nur begrenzt auf vielf\"altige Dokumenttypen verallgemeinern. Neuere generative Ans\"atze mit gro{\ss}en Sprachmodellen erzielen gute Ergebnisse bei der Textextraktion, doch die r\"aumliche Lokalisierung kann weniger genau sein, und generierte Ausgaben sind nicht auf beobachtete Eingabetoken beschr\"ankt.

\textbf{Methode:} Die Zwei-Kopf-Architektur weist jedem Token eine Rolle (Schl\"ussel, Wert oder Sonstiges) zu und verbindet vorhergesagte Schl\"ussel-Spans mit Wert-Spans. Wir untersuchen die Verkn\"upfung auf Token-Ebene (V1), auf Span-Ebene (V2), validierungsbasierte Diagnosen und Korrekturen sowie eine abschlie{\ss}ende korrigierte Konfiguration (V4). Da Quelldateien nicht verf\"ugbar oder besch\"adigt waren oder keine nutzbare Textebene enthielten, umfasst der aufbereitete Entwicklungssatz 5.389 der 9.656 eindeutigen KVP10k-Trainingsseiten, die anschlie{\ss}end in Trainings- und Validierungsteilmengen aufgeteilt werden.

\textbf{Evaluation:} Wir folgen der ver\"offentlichten KVP10k-Benchmarkimplementierung mit typabh\"angigem Eins-zu-eins-Matching, NED $<0{,}2$, IoU $\geq0{,}3$ im ausf\"uhrbaren Code sowie Makromittelung von Pr\"azision und Recall pro Seite. Entit\"atsextraktion und Schl\"ussel-Wert-Verkn\"upfung werden getrennt bewertet.

\textbf{Ergebnisse:} Die V4-Checkpoint-Auswahl verwendet den makrogemittelten F1-Wert f\"ur regul\"are Paare unter gemeinsamer Text- und Lagebewertung und w\"ahlt Epoche 10. Auf dem Testsatz erreicht V4 f\"ur regul\"are Paare F1-Werte von 0,392 (Text), 0,446 (Lage) und 0,345 (Text und Lage). Die korrigierten klassenbewussten Mikro- und Makro-F1-Werte der Entit\"atsklassifikation betragen 0,792 und 0,772. Die rekonstruierte Mistral-Baseline erreicht 0,521. Eine unver\"anderte Nachverarbeitung erh\"oht den gemeinsamen F1-Wert \"uber alle Kategorien von 0,229 auf 0,261, ohne den Wert f\"ur regul\"are Paare zu ver\"andern. Die numerisch h\"oheren Entit\"atsergebnisse zusammen mit den qualitativen Fehlerbeispielen legen nahe, dass die Relationsverkn\"upfung weiterhin eine wesentliche Einschr\"ankung der aktuellen Pipeline darstellt.

\textbf{Stichw\"orter:} Dokumentenverstehen, Schl\"ussel-Wert-Paar-Extraktion, LayoutLMv3, Span-Ebenen-Verkn\"upfung, Layoutbewusstes Lernen, Informationsextraktion
}

\clearpage
